\documentclass[../../../include/open-logic-section]{subfiles}

\begin{document}

\olfileid{sth}{choice}{tarskiscott}
\olsection{Tarski--Scott 技巧}

在\olref[cardinals][cardsasords]{defcardinalasordinal}中，我们将基数定义为序数。要做到这一点，我们假定了良序公理。我们这样做纯粹是因为它是“行业标准”。

因此，在我们讨论与良序公理相关的任何哲学问题之前，必须明确一点：我们\emph{可以}偏离行业标准，并且\emph{无需}假定良序公理来发展基数理论。我们仍然可以采用 $\cardeq{A}{B}$、$\cardle{A}{B}$ 和 $\cardless{A}{B}$ 的定义，即\olref[sfr][siz][]{chap}中所给出的那些定义。我们只需要一个新的\emph{基数}概念。

一种朴素的想法是试图这样定义 $A$ 的基数：\[
	\Setabs{x}{\cardeq{A}{x}}.
\]
你也许想把它与 Frege（弗雷格）对 $\# x Fx$ 的定义作个比较，这在\olref[cardinals][hp]{sec}末尾曾有概述。而且，由于我们在那里提到的原因，这个定义行不通。任何单元素集都与 $\{\emptyset\}$ 等势。但新的单元素集在层级的每个后继阶段都会产生（只需考虑前一阶段的单元素集即可）。因此，$\Setabs{x}{\cardeq{A}{x}}$不存在，因为它不可能具有秩。

为了解决这个问题，我们使用一个由 Tarski（塔斯基）和 Scott（斯科特）提出的技巧：\footnote{提醒：所有公式都可能含有参数（除非另有明确说明）。}

\begin{defn}[Tarski--Scott]
对于任意公式 $\phi(x)$，令
$[ x : \phi(x)] $ 为所有满足 $\phi(x)$ 且具有最小可能秩的 $x$ 的集合（或者，如果没有 $\phi$ 这样的 $x$，则为 $\emptyset$）。
\end{defn}

我们应当验证这个定义是否合法。在 $\ZF$ 中，\olref[spine][foundation]{zfentailsregularity} 保证了每个 $x$ 的 $\setrank{x}$ 都存在。现在，如果存在任何满足 $\phi$ 的对象，那么我们可以令 $\alpha$ 为满足 $(\exists x\subseteq V_\alpha)\phi(x)$ 的最小秩，即 $(\forall \beta
\in \alpha)(\forall x \subseteq V_\beta)\lnot \phi(x)$。然后我们可以通过分离公理将 $[x : \phi(x)]$ 定义为 $\Setabs{x \in
V_{\alpha+1}}{\phi(x)}$。

在论证了 Tarski--Scott 技巧的合理性之后，我们现在可以用它来定义基数概念：

\begin{defn}
$A$ 的 \textsc{ts}-基数是 $\text{tsc}(A) = [x :
\cardeq{A}{x}]$。
\end{defn}

\textsc{ts}-基数的定义没有使用良序公理。但是，即使没有该公理，我们也能证明\emph{\textsc{ts}-基数}的表现与\olref[cardinals][cardsasords]{defcardinalasordinal}中定义的\emph{基数}非常相似。例如，如果我们用 \textsc{ts}-基数来重述 \olref[cardinals][cardsasords]{lem:CardinalsBehaveRight} 和 \olref[card-arithmetic][opps]{lem:SizePowerset2Exp}，那么证明在 $\ZF$ 中完全行得通，无需假定良序公理。

顺便一提，我们也可以使用 Tarski--Scott 技巧来发展序数理论。对于良序 $\tuple{A, <}$，令 $\text{tso}(A, <) = [\tuple{X,
R} : \ordeq{\tuple{A, <}}{\tuple{X, R}}]$。有关基数和序数这种处理方式的更多内容，请参见 \citet[chs.~9--12]{Potter2004}。

\end{document}